\section{A fixed-floor Fourier and resource transfer}\label{ut:section}

This appendix proves the interpretation of the fixed floor-$25$
acquisition used in \eqref{loc:weighted-transfer}. The floor is an
evaluation parameter in a collection of inequalities. It is not another
sample size at which the current law is assumed to be extremal.
The moment-defect identity and the underlying enclosure method retain
their attribution to Rodziewicz--Ko\l{}odziejek~\cite{RK2026}.

Fix $N=25$, an integer $n\in\{26,\ldots,30\}$, and a standardized law
$P$ attaining the unrestricted $C_n=C\ge\CE$ at a positive closed
threshold $H$. Assume $C_{n-1}\le\CE$ and
$\rho\in[1,11/10]$ or $\rho\in[11/10,6/5]$.
Use $D$, $\mathcal L$ and $\mathcal R_n$ from
\eqref{loc:conserved}--\eqref{loc:raw}, all for this actual pair $(P,H)$.
Write $\mathcal R_n=\mathcal R$ for the right side of
the conserved identity \eqref{loc:conserved} at this same law.
We require an actual-index score cap $C\le c_{\rm in}$ with
$c_{\rm in}\le434477/10^6$; the Gaussian replacement bound
$C<.422$ under the sharp predecessor is sufficient.
The support and moments used below must have been acquired for this
same actual-$n$ pair. No assertion about $C_{25}(P)$ is required.

\subsection{The complete low-frequency estimate}

Put $f(t)=\E e^{itX}$, $\gamma=\E X^3$, and
\begin{align*}
 V_n(t)&=f(t)^n-e^{-nt^2/2}(1-in\gamma t^3/6),\\
 I_0(n,H)&=\int_0^{1/2}\frac{|\Im(e^{-itH}V_n(t))|}{t^3}\,dt,\\
 I_\infty(n)&=\int_{1/2}^{\infty}\frac{|V_n(t)|}{t^3}\,dt.
\end{align*}
By \eqref{loc:fourier-remainder},
\begin{equation}\label{ut:split}
 |E_n|\le\frac2\pi\{\sqrt n I_0(n,H)+\sqrt n I_\infty(n)\}.
\end{equation}
Both integrals are retained, including the entire Gaussian comparator
on $[1/2,\infty)$.

\begin{lemma}[Low-frequency candidates]\label{ut:low}
Suppose $n>2$, $M_4\le M\le111/10$, $|\gamma|\le g\le37/20$,
and $M_6\le\overline M_6$. For the actual coefficients
\[
 a=(3-M_4)/24,\qquad z_5=(10\gamma-\E X^5)/120,
 \qquad c_5=a\gamma/6,
\]
let $a_*\ge|a|$ and $Z\ge|z_5|$. Define, for real $j>2$,
\begin{gather*}
 A_j=\frac{j-1}{2j},\quad B_j=\frac{45-M}{96}(1-2/j),
 \quad F_j=(j/(j-1))^{3/2},\quad I_u\ge\sqrt{\pi/2},\\
 E_6=\frac{\overline M_6}{720}+\frac{g^2}{72}
                  +\frac{(1/2+g^2/144)^3}{5},\quad
 T_6=a_*+Z/2+E_6/4,\quad
 L_2(j)=\frac{E_6}{2A_j^2}+\frac{T_6^2}{2B_j^3}.
\end{gather*}
Each of the following complete triples satisfies
\begin{equation}\label{ut:low-bound}
 \sqrt n I_0(n,H)\le L_0(n)+L_H(n)|H|+L_2(n)/\sqrt n.
\end{equation}
The triangle triple is
\begin{align}\label{ut:triangle}
 L_0(j)&=I_uF_j(Z+a_*g/2)+251g^3/17280,\notag\\
 L_H(j)&=I_uF_j a_*+251g^2/4800.
\end{align}
For the cancellation triple, let a closed rectangle contain the actual
$(z_5,c_5)$, and let $J_*\ge0$ bound
\[
 J(z,c)=\int_0^\infty e^{-x^2/2}|zx^2-cx^4|\,dx
\]
at every corner of that rectangle. Then one may take
\begin{align}\label{ut:paired}
 L_0(j)&=F_jJ_*+\frac{F_jI_u}{j-1}
       (35a_*g^3/48+35Zg^2/24)+251g^3/17280,\notag\\
 L_H(j)&=F_jI_u\{a_*+5Zg/[2(j-1)]\}+251g^2/4800.
\end{align}
In either case, all three coefficients and $L_2(j)/\sqrt j$ are
nonincreasing for $j\ge N$ when their physical inputs are fixed.
\end{lemma}
\begin{proof}
Set $h(t)=e^{-t^2/2-i\gamma t^3/6}$. Taylor expansion gives
\[
 f(t)-h(t)=-at^4-iz_5t^5+r_6(t),\qquad
 |r_6(t)|\le E_6t^6\quad(0\le t\le1/2).
\]
Indeed, for $u=-t^2/2-i\gamma t^3/6$,
$|u|\le t^2(1/2+g^2/144)<1/6$;
the omitted exponential terms cost at most $|u|^3/5$, and its
degree-six quadratic term costs $g^2t^6/72$.
Symmetrization and $\cos y\le1-y^2/2+y^4/24$ give
\[
 |f(t)|,|h(t)|\le e^{-(45-M)t^2/96}.
\]
Taylor expansion of the polynomial $w^n$ along the joining segment
therefore yields
\[
 f(t)^n-h(t)^n=n(f(t)-h(t))h(t)^{n-1}+R_2(t),\qquad
 |R_2(t)|\le\tfrac12n^2T_6^2t^8e^{-nB_nt^2}.
\]
Integrating $nr_6h^{n-1}$ and $R_2$, after division by $t^3$,
costs at most $L_2(n)/n$.
For $u=n\gamma t^3/6$, rotation by $Ht$ bounds the remaining
Gaussian comparator's imaginary part by
$e^{-nt^2/2}(|u|^3/6+|H|t u^2/2)$.
The Gaussian even moments give its integral at most
$n^{-1/2}(251g^3/17280+251g^2|H|/4800)$.

For the leading quartic and fifth terms put $b=n-1$ and
$v=b\gamma t^3/6$. The bounds
$|\sin(Ht+v)|\le|H|t+bg t^3/6$ and $|\cos|\le1$
give \eqref{ut:triangle}.
Alternatively, the leading imaginary part after division by $t^3$
is bounded, before multiplication by $ne^{-bt^2/2}$, by
\[
 |z_5t^2-atv|+a_*t|\sin v-v|+Zt^2|\cos v-1|
                 +|H|t\{a_*t+Zt^2|\sin v|\}.
\]
The first integral is $nJ(z_5,c_5)/b^{3/2}$.
The norm $J$ is convex in $(z,c)$, so its rectangle maximum is
bounded by the largest of its corner bounds. Integrating the other
four terms with the global sine and cosine remainders gives
\eqref{ut:paired}. These arguments use the original $H$, without
rescaling it to a floor-$N$ threshold.

Finally,
\[
 \frac{F_j'}{F_j}=-\frac3{2j(j-1)},\quad
 A_j'=\frac1{2j^2}>0,\quad B_j'=\frac{45-M}{48j^2}>0.
\]
Thus $F_j$, $F_j/(j-1)$, $A_j^{-2}$, $B_j^{-3}$ and
$j^{-1/2}$ decrease. Every multiplier of these factors in
\eqref{ut:triangle}--\eqref{ut:paired} is nonnegative, including
$J_*$. This proves each claimed coefficient comparison.
\end{proof}

Here are the coefficient choices actually needed in these candidates.
They are physical moment inequalities, independent of the sample size.
For the triangle choice, assume $|X|\le R\le6$ and $M_5\le\eta$.
One may use $a_* =\max(2,M-3)/24$, $\overline M_6=R\eta$, and
$120Z=\min\{\eta+10g,Z_6,Z_r\}$, where
\[
 Z_6=\begin{cases}
 6(M-1)+(10-2M)g-6g^2+g^3,
       &M\le5,\ 10-2M-12g+3g^2\ge0,\\
 6(M-1)+(10-2M)^2/[4(6-g)],&M\le5\text{ otherwise},\\
 6(M-1)+(2M-10)^2/24,&M>5,
 \end{cases}
\]
and, with $r_0=\max(R,2g)>g$, put
\[
 t_*=\begin{cases}
 \min(g,(10-2M)/(2(r_0-g))),&M\le5,\\
 \min(g,(2M-10)/(2r_0)),&M>5.
 \end{cases}
\]
Then
\[
 Z_r=\begin{cases}
 r_0(M-1)+(10-2M)t_*-(r_0-g)t_*^2,&M\le5,\\
 r_0(M-1)+(2M-10)t_*-r_0t_*^2,&M>5.
 \end{cases}
\]
To verify them, reflect only the auxiliary Fourier integral so that
$t=\gamma=|\gamma|\ge0$, and set $Q=X^2-tX-1$.
Direct multiplication gives
\begin{equation}\label{ut:pearson}
\begin{aligned}
 \E Q^2&=M_4-1-t^2,\\
 \E X^5&=t^3+2t+\E[Q^2(X+2t)],\\
 M_6&=1+3t^2+t^4+\E[Q^2\{(X+t)^2+2t^2+2\}].
\end{aligned}
\end{equation}
For an enclosing radius $r\ge2g$, the two signs of
$\E X^5-10t$ are bounded by
$r(M-1)\pm(10-2M)t-rt^2\pm t^3$.
Use $t^3\le gt^2$, or discard the negative cubic, and maximize the
resulting concave quadratic on $[0,g]$. At $r=6$, the increasing
endpoint alternative gives the first line of $Z_6$.

The resource implementation also uses the following general signed
rectangle, which does not require $M<3$ or $R>2g$.
Suppose $0\le g_l\le|\gamma|\le g$, $M\ge1+g_l^2$, and $R\ge1$.
From \eqref{ut:pearson}, for every feasible $t\in[g_l,g]$,
\begin{gather*}
 L(t)\le120z_5\le U(t),\\
 L(t)=-R(M-1)+(10-2M)t+Rt^2+t^3,\\
 U(t)=\begin{cases}
 R(M-1)+(10-2M)t-Rt^2+t^3,&t\le R/2,\\
 8t-t^3,&t\ge R/2.
 \end{cases}
\end{gather*}
The second line for $U$ uses $\E Q^2\ge0$ when $R-2t<0$;
using $M-1-t^2$ there would have the wrong direction.
Choose outward bounds $L_*\le\min L$ and $U_*\ge\max U$ on the
indicated closed intervals, and set
\begin{gather*}
 z_l=\max(-\eta-10g,L_*)/120,\quad
 z_u=\min(\eta+10g,U_*)/120,\quad Z=\max(|z_l|,|z_u|),\\
 a_* =\max(|3-M|,|2-g_l^2|)/24,\\
 c_l=\frac{3-M}{144}
       \begin{cases}g_l,&M\le3,\\g,&M>3,\end{cases}
 \qquad c_u=\max_{g_l\le t\le g}(2t-t^3)/144,\\
 \overline M_6=\min\{R\eta,
  1+3g^2+g^4+(R^2+2Rg+3g^2+2)(M-1-g_l^2)\}.
\end{gather*}
Check $z_l\le z_u$, $c_l\le c_u$ before using the rectangle.
The cubic extrema can be enclosed by the Bernstein identity on a
finite closed subdivision; this is an enclosure of the entire
interval, rather than evaluation at finitely many sample points.
The displayed $a_*$ uses $1+g_l^2\le M_4\le M$.
The last line follows from the final identity in \eqref{ut:pearson},
with all substituted factors nonnegative. These prove the hypotheses
of Lemma~\ref{ut:low} for the general signed candidate used by both
resource branches. The restricted $M<3$, $R>2g$ formula is not needed
to license a call outside its own domain.

The resource branch enables this signed candidate only on an already
acquired interval $0<l\le p\le u<1/2$, with $\theta\le\Theta<1$
and $A_-\ge A_l>0$, $A_+\ge B_l>0$.
Write $\rho_B(p)=\{1-2p(1-p)\}/\sqrt{p(1-p)}$ and set
\begin{gather*}
 \rho_* =\max\left\{1,\rho_l,\rho_B(u)
       -\max(0,3\rho_B(l)/2-3\min(A_l,B_l))\Theta\right\},\\
 g_l=\max\left\{0,\sigma_l\rho_*-3-d_D,
 \frac{(1-\Theta)^{3/2}(1-2u)}{\sqrt{u(1-u)}}
                -(R+2\sqrt{u/(1-u)})\Theta\right\}.
\end{gather*}
Here $D\le d_D$ is an incoming bound. These are genuine physical
lower bounds for $\rho$ and $|\gamma|$, respectively.
For the first, conditional expansion gives
$\rho\ge(1-\theta)^{3/2}\rho_B(p)+3\min(A_l,B_l)\theta$;
use $(1-\theta)^{3/2}\ge1-3\theta/2$ and the monotonicity of
$\rho_B$ on $(0,1/2)$.
For the second, $\gamma\ge\sqrt{10}\rho-3-D$ follows from
$J\le\rho+1$, whereas conditional expansion about the two means gives
\[
 \gamma\ge\frac{(1-\Theta)^{3/2}(1-2u)}{\sqrt{u(1-u)}}
                     -(R+2\sqrt{u/(1-u)})\Theta.
\]
In the latter calculation the positive-class residual is nonnegative
and the negative-class residual has magnitude at most the displayed
coefficient times $\theta$.
The gates are $\rho_*\le\rho_u$, $0\le g_l\le g\le37/20$,
$1\le M\le111/10$, $R\ge1$ and $M-1-g_l^2\ge0$.
The same $\rho_*$ is kept in the paired rectangle, score corners,
phase callback and exports. On another probability half, or if this
candidate is ineligible, the unpaired candidate remains available;
ineligibility alone does not exclude a law.

For completeness, the smaller Gaussian corner prices used in this
rectangle also have an index-free proof. If $zc\le0$, then
$J(z,c)=\sqrt{\pi/2}(|z|+3|c|)$.
If $zc>0$, replace $(z,c)$ simultaneously by $(|z|,|c|)$ and put
$x=z/c>0$. Splitting the integral at $\sqrt x$ gives
\begin{equation}\label{ut:mills-norm}
\begin{aligned}
 J(z,c)&=\sqrt{\pi/2}(z-3c)+\sqrt{zc}e^{-x/2}
       \left(6+\frac{2(3-x)}x\sqrt x\,\mathcal M(\sqrt x)\right),\\
 \mathcal M(r)&=\int_0^\infty e^{-rt-t^2/2}dt.
\end{aligned}
\end{equation}
For $I_j(r)=\int_0^\infty t^je^{-rt-t^2/2}dt$, integration by parts
gives $I_0=1/(r+I_1/I_0)$ and
$I_j/I_{j-1}=j/(r+I_{j+1}/I_j)$.
Truncating this finite recursion with a zero last ratio gives an
alternating upper bound at even depth and lower bound at odd depth.
In rational form, set $d_{k+1}=1$,
$d_j=1+j/(xd_{j+1})$ for $j=k,\ldots,1$.
In \eqref{ut:mills-norm} replace
$\sqrt x\mathcal M(\sqrt x)$ by $d_1^{-1}$ with $k=11$ if
$x\ge3$, and $k=12$ if $x<3$.
The sign of $3-x$ gives the required upper bound in both cases.
The remaining multiplier is positive; for $x\ge3$ it exceeds
$6-2(x-3)/x>4$.
A negative coefficient $z-3c$ uses a lower enclosure of
$\sqrt{\pi/2}$ and a positive coefficient uses an upper enclosure.
The minimum with
\[
 \sqrt{\pi/2}|z-3c|+6\sqrt{zc}e^{-x/2}
\]
is also a valid corner bound. These are the two corner alternatives
in the saved candidates. Auxiliary reflection here leaves the original
signed $\gamma$ in the contact identity unchanged.

\subsection{The complete high-frequency estimate}

Let $X'$ be an independent copy of this same $X$. A candidate consists
of fixed numbers $M,a,m,d,w$ and an actual event $E$ such that
\begin{equation}\label{ut:event-input}
\begin{gathered}
 M_4\le M\le111/10,\qquad\PP(E)\ge a>0,\\
 |X-X'|\ge d>0\text{ on }E,\qquad
 \E[(X-X')^2\mid E]\le m,\\
 0<d\le12,\qquad0<w\le1/4.
\end{gathered}
\end{equation}
Here $a\le1$ and $m>0$. Put $b=9/10$,
\begin{gather*}
 h=333/106-(22/7)^3w^2/6>0,\quad
 \kappa_E=a(23/48)^2d^2,\quad
 \Lambda(x)=\frac{x+2}{x^2+4x+2},\\
 S_3=\sum_{j=1}^{20}j^{-3}+1/800,\quad
 S_5=\sum_{j=1}^{20}(j-w)^{-5}+[4(20-w)^4]^{-1},\\
 c=\tfrac12\left(1-\max_{u\in\{1/4,81/100\}}
                    [1-u+(M+3)u^2/12]\right),\\
 c_0=1/8-(M+3)/384,\qquad s=\{1-(M+3)53/600\}/2.
\end{gather*}
Define
\begin{align}\label{ut:high-six}
 \mathfrak H_j={}&A_j^{\rm high}
 +\frac{m}{4\pi^{3/2}h\sqrt a}
                \left(S_3+\frac{3S_5}{jah^2}\right)
 +\frac{m\sqrt j}{2\pi^2}e^{-jah^2w^2}\notag\\
 &+Z_j+\frac{\sqrt j}{2b^2}e^{-ja(23/48)^2}
 +\sqrt j e^{-j/8}\{2\Lambda(j/8)+\sqrt{M-1}/3\}.
\end{align}
One always may choose
\[
 A_j^{\rm high}=2\sqrt j e^{-jc},\qquad
 Z_j=\frac{e^{-j\kappa_Eb^2}}{2b^4\kappa_E\sqrt j}.
\]
If $s>0$, the first may be replaced by its minimum with
$2\sqrt j e^{-jc_0}\Lambda(js/4)$.
If $j\kappa_Eb^2\ge1/4$, the second may be replaced by its minimum
with $\sqrt j e^{-j\kappa_Eb^2}\Lambda(j\kappa_Eb^2)/(2b^2)$.

\begin{lemma}[High-frequency candidates and their gates]\label{ut:high}
Under \eqref{ut:event-input}, $\sqrt n I_\infty(n)\le\mathfrak H_n$.
A fixed choice of alternatives satisfies
$\mathfrak H_n\le\mathfrak H_N$ for every $n\ge N$ provided
\begin{equation}\label{ut:high-gates}
 2Nc'>1\quad
 (c'\in\{c,ah^2w^2,a(23/48)^2,1/8\}).
\end{equation}
For the improved first term also require $s>0$ and $2Nc_0>1$;
for the improved zero-cell term require $N\kappa_Eb^2\ge1/4$.
An alternative failing its floor gate is omitted, even if it would
become eligible at a larger index.
\end{lemma}
\begin{proof}
On $[1/2,b]$, symmetrization gives
$|f(t)|^2\le1-t^2+(M+3)t^4/12$, proving the first term with rate
$c$. The sharper chord is
\[
 t^2/2-(M+3)t^4/24\ge c_0+s(t^2-1/4).
\]
It is valid since the difference is
$(M+3)(t^2-1/4)(81/100-t^2)/24\ge0$.

For $t\ge b$, let $p_E=\PP(E)$. Symmetrization, $1-u\le e^{-u}$,
and conditional Jensen imply
\[
 |f(t)|^n\le e^{-n\E\sin^2(t(X-X')/2)}
 \le\E[e^{-na\sin^2(t|X-X'|/2)}\mid E].
\]
For each conditional distance $D_*=|X-X'|$, partition the positive
axis into the zero cell and the cells centred at
$2\pi j/D_*$, $j\ge1$. In a nonzero cell write
$t=(2\pi/D_*)(j+v)$. On $|v|\le w$,
$|\sin(\pi v)|\ge h|v|$; on the wings it is at least $hw$.
The paired Taylor bound
\[
 \{(j-v)^{-3}+(j+v)^{-3}\}/2
             \le j^{-3}+6(j-w)^{-5}v^2
\]
and the Gaussian zeroth and second moments prove the second term
of \eqref{ut:high-six}. Its Jacobian is
$D_*^2/(4\pi^2)$, and conditional averaging uses
$\E[D_*^2\mid E]\le m$. Summing the two positive series and
bounding their tails after $20$ by integrals gives $S_3,S_5$.
The wing weight is at most $D_*^2/(2\pi^2)$ and gives the third term.
All extensions of these integrals below $b$ add nonnegative quantities.

In the zero cell, $u=tD_*/2\le\pi/2$.
For $u\le1/2$ use $\sin u\ge(23/24)u$ and $D_*\ge d$.
This gives the zero-cell Gaussian tail $Z_j$.
For $u\ge1/2$ use $\sin u\ge23/48$ and total weight
$1/(2b^2)$, giving the fifth term.
The comparator on the whole interval $[1/2,\infty)$ costs
\[
 \sqrt j\int_{1/2}^\infty e^{-jt^2/2}t^{-3}dt
 +\frac{j^{3/2}|\gamma|}{6}\int_{1/2}^\infty e^{-jt^2/2}dt.
\]
The second integral is at most $2e^{-j/8}/j$, and Pearson gives
$|\gamma|\le\sqrt{M-1}$. This proves the sixth term.

The rational Gaussian-tail estimate used above is, for $x=jk b^2>0$,
\begin{equation}\label{ut:rational-tail}
 \sqrt j\int_b^\infty e^{-jkt^2}t^{-3}dt
 \le\frac{\sqrt j e^{-x}}{2b^2}\Lambda(x).
\end{equation}
After substitution its remaining factor is
$\int_0^\infty u(x+u)^{-1}e^{-u}du$.
For $P(u)=u^2-4u+2$ and $D_x=x^2+4x+2=P(-x)$,
division of $D_x^2-P(u)^2$ by $x+u$, followed by
$\int u^ke^{-u}du=k!$, gives
$\int_0^\infty (x+u)^{-1}e^{-u}du\ge(x+3)/D_x$.
The discarded remainder is $P(u)^2/[(x+u)D_x^2]\ge0$.
Consequently the preceding factor is at most
$1-x(x+3)/D_x=\Lambda(x)$, proving
\eqref{ut:rational-tail}. The simpler zero-cell alternative follows
by replacing $t^{-4}$ by $b^{-4}$ after integrating once.

For the floor comparison,
$(\sqrt j e^{-cj})'=(\sqrt j e^{-cj})(1/(2j)-c)$,
and $j^{-1/2}e^{-cj}$ decreases for every $c>0$.
Also $\Lambda'(x)=-(x^2+4x+6)/(x^2+4x+2)^2<0$.
These prove the assertions for the second, third, fifth and sixth
terms and for either first-term choice under the stated gates.
For the improved zero-cell choice the derivative of
$\sqrt x e^{-x}\Lambda(x)$ has, after multiplication by positive
factors, numerator
\[
 4-10x-22x^2-13x^3-2x^4.
\]
It equals $-11/128$ at $x=1/4$ and decreases for $x>0$.
This proves precisely its additional gate.
\end{proof}

The event inputs transfer without an index assumption. For $0<d^2<2$
and $E=\{(X-X')^2\ge d^2\}$, Cantelli's inequality supplies
\[
 V=2M+2,\quad h_0=2-d^2>0,\quad
 a=h_0^2/(V+h_0^2),\quad m=2+V/h_0.
\]
Indeed $Y=(X-X')^2$ has $\E Y=2$ and
$\operatorname{Var}Y\le V$. Applying the elementary one-sided
second-moment bound to $Y-2$ gives $\PP(Y\ge d^2)\ge a$.
For $p_E=\PP(E)$, Cauchy--Schwarz applied to the covariance gives
$\E[Y\mid E]\le2+\sqrt{V(1-p_E)/p_E}\le2+V/h_0$.
This also licenses the older weaker parameters
$a_U=h_0^2/(V+4)$, $m_U=2(V+4)/h_0^2$:
$a\ge a_U$ and
\[
 m_U-m=\frac{(2-h_0)(V+2h_0+4)}{h_0^2}\ge0.
\]
Every event-dependent term of \eqref{ut:high-six} decreases in $a$
and increases in $m$; its gates are preserved when $a$ increases.
Thus no separate universal-event proof is needed, and old certificates
using the weaker parameters remain valid. This domination holds for
the same $M,d,w$; it does not identify different distance choices.
The selected Cantelli pool uses
$d\in\{1/2,11/20,3/5,13/20,7/10\}$ and $w=1/4$.
Only members satisfying \eqref{ut:high-gates} at $N$ are retained.
Replacing $M$ by a fixed outward ceiling on the $1/1000$ grid
is legitimate as a new fourth-moment cap; it is then held fixed.

After a central gap $(-A_-,A_+)$ is proved, another actual event is
$E=\{XX'<0\}$. With $p=\PP(X>0)$, $\mu=\E X_+$ and
$\beta=\E[X|X|]$, its exact conditional distance is
\[
 \PP(E)=2p(1-p),\qquad
 \E[(X-X')^2\mid E]
 =\frac{1+4\mu^2+\beta}{2p}
  +\frac{1+4\mu^2-\beta}{2(1-p)}.
\]
The expression increases in $\mu^2$, is affine in $\beta$, and is
convex in $p$ for $|\beta|\le1$. Incoming interval endpoints
therefore give $m$; the independent bound $m\le2/a$ may be intersected.
An acquired $D\le d_D<(\sigma_l-1)\rho_l$, where
$1<\sigma_l\le\sqrt{10}$, licenses the central gap. Moreover
\begin{equation}\label{ut:gap-input}
\begin{gathered}
 (A_-+A_+)^2\ge\frac{\rho_l-d_D/(\sigma_l-1)}{\mu},\qquad
 (A_-+A_+)^2\ge\frac{2P_*}{3\mu},\\
 P_* =\max\{(\sigma_l-1)\rho_l,\sigma_l\rho_l-\gamma_u\}-d_D.
\end{gathered}
\end{equation}
The first follows from the gap decomposition's
$T_++T_-\le D/(\sqrt{10}-1)$.
The second follows from
$J\le\rho+2\mu A_-A_+$ and hence $6\mu A_-A_+\ge P_*$.
Use an incoming upper bound for $\mu$ and positive lower roots.
The weaker first bound with $\mu\le1/2$, and independently acquired
$A_-+A_+\ge A_l+B_l$, give the other saved distance choices.
The gates are a positive gap, positive mean denominator and $0<d\le12$;
no assumption $d_D<2$ is needed for \eqref{ut:gap-input}.

A finite minimum of complete high-frequency candidates is valid:
each bounds the same $\sqrt n I_\infty(n)$.
Likewise a finite minimum of complete low-frequency triples is valid.
One may not combine the constant of one low triple with the slope of
another without a separate proof. Rational outward evaluations at
$N$ bound the exact displayed expressions at $N$ and hence at $n$;
no monotonicity of the rounding or exponential evaluator is assumed.

\subsection{Raw and charged candidates at the fixed floor}

Choose one complete low triple and one complete high bound at $N$.
For a fixed rational $\omega_\pi\ge2/\pi$, put
\begin{equation}\label{ut:fourier-floor}
 K_F=\omega_\pi\{L_0(N)+\mathfrak H_N+L_2(N)/\sqrt N\},
 \qquad s_F=\omega_\pi L_H(N).
\end{equation}
Lemmas~\ref{ut:low} and \ref{ut:high} prove
$|E_n|\le K_F+s_F|H|$ for the actual law.
The values $16/25$ and $212/333$ used in the two enclosure families
both satisfy the normalization condition.

Here is a precise description of the coefficient estimates which can
be used with \eqref{ut:fourier-floor}. On an incoming physical cell
assume that the coefficient atlas has proved
\begin{equation}\label{ut:physical-coefficients}
 3\beta J/2+3\rho J_1-3J_2-K_q\le A(\rho),\qquad
 K_1\in[k_l,k_u],\quad U\in[u_l,u_u].
\end{equation}
These are inequalities in the moments and support of $P$; they contain
no sample index. For the initial cover they are proved on the signed
$(\rho,\beta,\gamma,\mu)$ rectangles. For the resource cover they use
the conditional moment and gap formulas of Section~\ref{loc:section}.
The full physical derivations are the sections \emph{Signed moments
with a common first moment}, \emph{Preserving central and outer
integral costs} and \emph{Moment and gap coefficients for a whole
cell} in the continuous companion's \emph{Coefficient atlas for the
enclosure}, located in Appendix~\ref{app:companions}.
The central polynomial is $\rho P_1(A_-,A_+)-P_2(A_-,A_+)$, where
\begin{align*}
 P_1&=B^3/3+AB^2/2-A^3/6,\\
 P_2&=B^4/4+AB^3/3-A^4/12,
\end{align*}
where $(A,B)=(A_-,A_+)$.
Its full-rectangle Bernstein bound is a maximum of affine functions
of $\rho$. The outer cost is bounded by
$3\{(\rho-B_l)_+^2T_++(R-\rho)^2T_-\}/8$.
These and the remaining affine moment terms give the convex
$A(\rho)$ used in the corner reduction. Every signed product with
a possibly negative factor retains the appropriate two endpoints.
Thus \eqref{ut:physical-coefficients}, rather than a file status, is
the required physical-coefficient premise of a candidate.

Suppose previous information gives
$\mathcal R_n\le r$, $nD\le Nd$ and $0<r<2N/5$.
Once removal of the Gaussian weight has been licensed, it follows that
\begin{equation}\label{ut:incoming-envelope}
 D\le d,\qquad L_N\le\varphi\le2/5,\quad
 L_N=\frac{997}{2500}(1-5r/(2N))>0.
\end{equation}
The density statement uses the nonnegative phase term in
\eqref{loc:conserved}, $\CE\rho>2/5$ and $n\ge N$.
For the unseeded first pass use $L_N=0$, without assuming $r$ or $d$.

For a sign $\varepsilon\in\{-1,1\}$ and a prospective score $c$, set
\[
 F_\varepsilon(c)=s_F+
 \max_{k\in\{k_l,k_u\},\,u\in\{u_l,u_u\},\,phi\in\{L_N,2/5\}}
                  \varepsilon(ck+\phi u).
\]
The maximum of the two signs may be used for an unsigned bound.
If $\gamma\ge g_0(\rho)$ is an incoming signed lower bound, with
$g_0$ affine in $\rho$ (including a fixed signed interval endpoint),
the correct phase curvature is
\begin{equation}\label{ut:signed-curvature}
 q(c,\rho)=c\rho+\tfrac12
           \min\{L_N g_0(\rho),(2/5)g_0(\rho)\}.
\end{equation}
In particular $g_0(\rho)=3.16\rho-3-d$ is valid since
$J\le\rho+1$. The alternative general-envelope curvature is
\[
 q_{\rm gen}(c,\rho)=c\rho+L_Ng_0(\rho)/2
                    -(2/5-L_N)(d-.16)_+/2.
\]
It is no larger than \eqref{ut:signed-curvature}, because
$(-g_0(\rho))_+\le(d-.16)_+$ for $\rho\ge1$.
Either curvature must be proved strictly positive before division.
Using $L_Ng_0/2$ alone when $g_0<0$ would not be licensed.

For the actual $H$ with sign $\varepsilon$, \eqref{loc:raw} now gives
\begin{equation}\label{ut:raw-candidate}
 \mathcal R_n\le K_F+C A(\rho)+F_\varepsilon(C)|H|
                                      -q(C,\rho)H^2.
\end{equation}
Thus a floor candidate is an envelope for the actual raw threshold.
It is not an estimate of a fictitious $\mathcal R_{25}(P)$.

The same incoming bounds give a conserved-phase price
\begin{equation}\label{ut:rebate}
 e_N(\rho)=\tfrac12(1-5r/(2N))(.4095\rho-d/15)>0.
\end{equation}
Indeed the difference between the weighted and unweighted terms in
\eqref{loc:conserved} is
$n(1-e^{-w})(\CE\rho-\phi(0)D/6)$, $w=H^2/(2n)$.
Use $1-e^{-w}\ge we^{-w}$, the density floor before multiplying by
$\phi(0)$, $D\le d$ and $\phi(0)<2/5$.
Consequently, with $\kappa=997/15000$,
\begin{equation}\label{ut:charged-candidate}
 \kappa nD+n(C-\CE)\rho
 \le K_F+C A(\rho)+F_\varepsilon(C)|H|
                       -\{q(C,\rho)+e_N(\rho)\}H^2.
\end{equation}
The first unseeded pass uses $e_N=0$.

The own-score interval is also conservative at $N$. If a physical
cell forces $D\ge d_l\ge0$ and $\rho\ge\rho_l>0$, then
\begin{equation}\label{ut:own-score}
 C\le .40974+
 \frac{r/N-L_Nd_l/6}{\rho_l}
\end{equation}
whenever the numerator is nonnegative. A negative numerator excludes
the cell. This follows from \eqref{loc:conserved}, since
$r/n-\varphi d_l/6\le r/N-L_Nd_l/6$.
The independent moment envelope and $c_{\rm in}$ may be intersected
with \eqref{ut:own-score}.

For each candidate freeze all physical coefficients before taking the
score and moment corners. For fixed $H$, the right side of
\eqref{ut:raw-candidate} is convex in $c$ and in $\rho$ separately:
$A$ is convex, $F_\varepsilon$ is a maximum of affine functions,
and $q$ is affine in $c$ and concave in $\rho$.
The same is true after adding the affine curvature $e_N$.
Hence the whole physical rectangle is bounded by its corners, first
pointwise in $H$, and then by completing each corner's square.
For $Q>0$, $\sup_{h\ge0}(A+Fh-Qh^2)=A+(F_+)^2/(4Q)$;
using $A+F^2/(4Q)$ is a weaker valid bound.

For a charged candidate split the score interval at $.40974$ and use
\begin{equation}\label{ut:excess-charge}
 n(C-\CE)\rho\ge N(C-.40974)_+\rho.
\end{equation}
On either side of this breakpoint the subtracted right side is
affine in each of $c,\rho$ separately, so the same corner argument
applies. Omitting this breakpoint would not follow from convexity.
If $R_N$ is the maximum raw corner price and $B_N$ the maximum
charged corner price, the resulting conclusions are
\begin{equation}\label{ut:stage-output}
 \mathcal R_n\le R_N,\qquad nD\le B_N/\kappa.
\end{equation}
They remain true after minima with other complete candidate bounds.
The phase roots of these same corner quadratics bound the original
$H$, and the pointwise charged quadratics give
$D\le\max_j(A_j+F_{j,\varepsilon}|H|-Q_jH^2)/(\kappa n)$.
Their normalization is therefore $15000/(997n)$.
For an unpaired export the implementation may instead use an already
proved $|H|\le H_*$ and the scalar weighted defect bound to write
\[
 \kappa nD\le\kappa Nd
 \le\kappa Nd+\epsilon(H_*^2-H^2),\qquad\epsilon>0.
\]
This is a valid pointwise envelope on the acquired phase interval.
Its artificial quadratic term asserts no additional physical curvature.

\subsection{The unseeded step and finite induction}

\begin{proposition}[Fixed-floor interpretation]\label{ut:transfer}
Consider the fixed finite B0 or B1 acquisition schedule at $N=25$.
At each Fourier step retain a complete candidate satisfying
Lemmas~\ref{ut:low} and \ref{ut:high} at the floor, its proved
physical coefficient bounds \eqref{ut:physical-coefficients}, and
the positive corner divisors above. At each physical step retain only
necessary inequalities for the current law and a complete closed
cover of its incoming domain. Every support or strip implication is
proved at the actual integer $n$, using its own incoming moment and
threshold conditions before its output is used.
Then every saved raw and defect output $(r_j,d_j)$ has the interpretation
\begin{equation}\label{ut:induction-invariant}
 \mathcal R_n\le r_j,\qquad nD\le Nd_j.
\end{equation}
The physical moment, endpoint, sign and threshold bounds exported
alongside it apply simultaneously to the same actual-$n$ law.
\end{proposition}
\begin{proof}
For the first pass there is no input $r_j$ or $d_j$.
Proposition~\ref{gc:initial-support} supplies the initial radii
$12/5$ on B0 and $13/5$ on B1 uniformly in the actual integer.
Before any defect feedback,
$M_4\le R\rho$ gives the respective ceilings $66/25$ and $78/25$;
$M_5\le RM_4$, $M_6\le RM_5$ and $|\gamma|\le\rho\le6/5$
license the triangle low bound and the universal or Cantelli high
bound. The initial signed-moment cover uses $\varphi\in[0,2/5]$,
the independent score interval and physical moment inequalities.
It assumes neither a central gap nor a final conditional moment cap.

The selected first raw cover maxima, with outward arithmetic, are
\[
 r_1^{\rm B0}=109743487/50000000,
 \qquad r_1^{\rm B1}=283144217/100000000.
\]
These are finite upper bounds on the raw prices in
\eqref{ut:raw-candidate}; their closed-domain arithmetic is part of
the companion calculation. Both are below $3<.4N\le nC\rho$.
It follows from \eqref{loc:conserved} that
\[
 \frac{\varphi}{6}(\gamma+3J)
            =C\rho+\sqrt n\mathcal L-\mathcal R_n/n>0.
\]
Consequently $\CE\rho-\phi(0)D/6>0$. Removing the Gaussian weight
in \eqref{loc:conserved} now lowers its left side and proves
\begin{equation}\label{ut:bootstrap}
 n(C-\CE)\rho+\kappa nD+n^{3/2}\mathcal L\le\mathcal R_n.
\end{equation}
Only at this point is the charged output from the same first cover
used. Its corner calculation, with $e_N=0$ and
\eqref{ut:excess-charge}, gives
$d_1=B_N/(\kappa N)$, with the saved outward values
\[
 d_1^{\rm B0}=129264343/10^8,\qquad
 d_1^{\rm B1}=16619657/10^7.
\]
Thus $nD\le Nd_1$, not merely $D\le d_1$.
Both $d_1<2< (\sqrt{10}-1)\rho_l$, so the next pass may use the
central gap. This order explains why calculating raw and charged
corner arrays in one routine is not a circular use of a defect cap.

Suppose \eqref{ut:induction-invariant} and the physical bounds have
been established at a later incoming stage. Since $n\ge N$, the
weaker physical statement $D\le d_j$ is available. All moment,
common-mean, gap and conditional-variance deductions can therefore
use $d_j$ in exactly their displayed physical directions.
For example, the additional p-local mean caps are
\[
 \mu^2\le p(1+\beta)/2,\quad
 \mu^2\le(1-p)(1-\beta)/2,\quad \mu^2\le p(1-p).
\]
They follow from signed-class Cauchy--Schwarz and retain the same
meaning on both probability halves. The endpoint-product distances
in \eqref{ut:gap-input} and all signed Pearson rectangles above
are likewise independent of $n$.

The physical inputs now license the next fixed Fourier candidates.
Their exact low and high expressions at the actual $n$ are bounded
by their floor-$N$ values, by the two lemmas. The floor density
interval \eqref{ut:incoming-envelope} contains the actual density;
the floor own-score interval \eqref{ut:own-score} contains the
actual score; and \eqref{ut:rebate} is a lower price for the actual
conserved phase. Equations~\eqref{ut:raw-candidate} and
\eqref{ut:charged-candidate} therefore hold pointwise in the actual
$H$. Their whole-cell corner maxima yield the next
$\mathcal R_n\le R_N$ and $nD\le B_N/\kappa$.
Writing the latter output as $d_{j+1}=B_N/(\kappa N)$ restores the
weighted invariant at every charged stage.

An excluded physical cell has an empty necessary interval or a
strictly negative upper bound for a nonnegative resource.
Any surviving law lies in at least one retained closed cell, so
maximizing its new bounds over all retained cells is valid.
Intersections with preceding bounds, finite minima of valid
candidates and outward rounding preserve the invariant.
This proves the assertion by induction over the prescribed finite
schedule, including repeated local INITIAL covers, the resource
subiterations and the final exports. The proof invokes neither
convergence of the iteration nor extremality at a smaller sample size.
\end{proof}

\subsection{Actual support and the interpretation of geometry}

The support hypotheses in Proposition~\ref{ut:transfer} use the
actual-index Gaussian majorant \eqref{loc:support-test}.
The original contact function $\mathcal D_n$ is zero on the full
support; the majorant need only be nonnegative there.
Write $s_n=\sqrt{n/(n-1)}$ and $e_n=n(s_n-1)-1/2$.
The possible extra allowance is
$n\rho s_n(C_{n-1}-C)_+=0$ under the assumed full predecessor.
A support ray or an excluded strip is obtained by proving the right
side strictly negative throughout that set, uniformly in all incoming
moment, score and threshold parameters. The Gaussian differences,
$s_n$ and $e_n$ in that calculation use the actual $n$.
For a bounded $H$ test, that bound on $H$ must already have been
proved by an earlier phase envelope. In particular the estimate
$e_n<.384/n$, proved in the retained derivation only for $n\ge36$,
is not used in these five actual-index requests.

The initial signed-moment interval follows from the support-free
inequality \eqref{ut:initial-signed-moment}; it requires no radius
or higher-moment feedback.

At the first support call the only moment restrictions are the
specified rho band and its consequences, together with the actual
score cap. An intermediate call consumes the moment intervals which
have already been proved by the preceding resource stage. Its new
radius licenses moment improvements only after the actual-index
strict contact test has passed. A final moment cap cannot license
the first support call, and a successful support calculation at $25$
cannot license a support or strip request at $26,\ldots,30$.
This is the one kind of acquisition step for which the floor-index
substitution is deliberately not made.

The distinction between weighted and physical quantities is also
necessary when exporting noise. From an acquired physical inequality
\[
 D\ge b+g\delta+h_NNs,\quad
 \delta=a_0+a_1,\quad b\ge0,\quad g,h_N>0,
\]
and $nD\le Nd_j$, put $r=N/n$ and $R_j=d_j-b$.
If $R_j>0$, $d_0=R_j/g$ and $L=R_j/h_N$ give
\begin{equation}\label{ut:export-plane}
 D\le rd_j,\qquad
 D\ge b+g\delta+rh_N\lambda_n,\qquad
 \frac{\delta}{rd_0}+\frac{\lambda_n}{L}
                 \le1-(n/N-1)b/R_j\le1,
 \quad\lambda_n=ns.
\end{equation}
This proves the units of the actual defect, inward intercept and noise
price. The joint intercept $L$ stays fixed. Without dividing by $R_j$,
the original inequality reads
$g\delta+rh_N\lambda_n\le rd_j-b$.
A negative right side excludes the cell, and a zero right side forces
$\delta=\lambda_n=0$. This also specifies the cases $R_j\le0$:
zero noise is retained only when the actual scaled right side is zero.

In the positive-residual case $R_j>0$, suppose physical support also gives $s\le B\delta$,
$B>0$, and $p(1-p)\ge v_l>0$. Retain first the full scaled
plane in \eqref{ut:export-plane} and put
\[
 c=1-(n/N-1)b/R_j,\qquad
 H_N=\left(\frac1L+\frac1{NBd_0}\right)^{-1}.
\]
The same scalar elimination as in \eqref{loc:harmonic-consequence},
now using $\delta\ge\lambda_n/(nB)$, $h=rh_N$,
$\overline D=rd_j$ and $nr=N$, gives
\begin{equation}\label{ut:scaled-harmonic}
 \lambda_n\left(\frac1L+\frac1{NBd_0}\right)\le c,
 \qquad \lambda_n\le cH_N.
\end{equation}
This is the stronger scalar export used by the actual-index adapter.
For a nonempty cell $c\ge0$, and $b\ge0$ gives $c\le1$;
with this same support constant $B$, $H_N$ is therefore also
a conservative scalar cap. At $c=0$ the scaled plane retains
$\delta=\lambda_n=0$.

If only the deliberately weakened, unscaled joint plane
$\delta/d_0+\lambda_n/L\le1$ is retained, it gives
\[
 \lambda_n\le\left(\frac1L+\frac1{nBd_0}\right)^{-1},\qquad
 \theta\le\left(1+\frac{nv_l}{L}
                         +\frac{v_l}{Bd_0}\right)^{-1}
 \le\left(1+\frac{Nv_l}{L}
                         +\frac{v_l}{Bd_0}\right)^{-1}.
\]
These weaker premises license the displayed actual-$n$
scalar expression and the transported variance fraction. They alone
do not license replacing that scalar expression by $H_N$.
The stronger cap in \eqref{ut:scaled-harmonic} is available when the
full scaled plane is retained. Both harmonic expressions are separate
scalar consequences; the original simultaneous intercept in
\eqref{ut:export-plane} remains $L$.
All further signed refinements may use the transported physical
$\theta$ and the actual support, followed by a new application of
the same joint-defect inequality. This completes the mathematical
transfer needed for the five actual-index acquisition schedules.

% Formula sources read for this appendix (all frozen source files are unchanged):
% source/analytic/UNIFORM25_CANTELLI_PRODUCTGAP_REVIEW.txt;
% source/analytic/UNIFORM25_BETA0_SOURCE_AND_COMPOSITION_REVIEW.txt;
% source/cutoff25_20261011/new_root/initial25_b1_exact_v1/
% RETAINED_ANALYTIC_FRAMEWORK.tex, initial34.py, kernel.py, high_window_core.py;
% the B0 initial25_b0_v2 counterparts; resource25_b0_v1 and resource25_b1_v1/
% candidate_charged_score_v2.py, contracted_kernel.py, general_paired_exact.py,
% corrected_phase_callback_extended.py, resource25_core.py.
